\documentclass[11pt]{article}

\usepackage[T1]{fontenc}
\usepackage{amsmath,amssymb,amsthm}
\usepackage[a4paper,margin=1in]{geometry}
\usepackage{booktabs}
\usepackage[hidelinks]{hyperref}
\hypersetup{
  pdftitle={A Classical-Mass-Action Sensitivity Zero in an Exact Six-Complex Census},
  pdfauthor={Xiaojian Zeng}
}

\newtheorem{theorem}{Theorem}
\newtheorem{lemma}{Lemma}
\newcommand{\R}{\mathbb{R}}
\newcommand{\Q}{\mathbb{Q}}
\newcommand{\diag}{\operatorname{diag}}
\newcommand{\rank}{\operatorname{rank}}

\title{A Classical-Mass-Action Sensitivity Zero in an Exact Six-Complex Census}
\author{Xiaojian Zeng\\
  \small PBC School of Finance, Tsinghua University, 100190\\
  \small \texttt{xiaojianzeng@foxmail.com}}
\date{10 September 2026}

\begin{document}
\maketitle

\begin{abstract}
For the ordered complexes
\(0,X,Y,2X,X+Y,2Y\), consider every simple directed reaction network
with three or four reactions, modulo exchange of the species, subject to
rank two, existence of a strictly positive steady flux, and a nonzero
independent-symbol Jacobian determinant.  We determine exactly whether
classical mass action can turn a generically nonzero inverse sensitivity
into an identically zero one.  The answer is affirmative.  In the prescribed
lexicographic order the minimum is
\[
  (0\to X,\ 0\to2X,\ X\to Y,\ X+Y\to0),
  \qquad (\rho,m)=(0\to X,Y).
\]
The cancellation is the steady-state identity
\(-v_{X\to Y}+v_{X+Y\to0}=-(Sv)_Y\), while the relevant determinant is
nonzero throughout the positive steady-flux locus.  Minimality is certified
by a symmetry-reduced enumeration using exact rational arithmetic and by a
separately implemented exact replay through the first witness.  The result is
only for this finite census; it is not a classification of mass-action
sensitivity zeros in general.
\end{abstract}

\section{Context and finite question}

Structural sensitivity analysis seeks steady-state responses that are
determined by a reaction network without fixing numerical rate laws.  The
augmented matrix used below arose in that program
\cite{mochizuki-fiedler-2015,brehm-fiedler-2018}.  Kuhrs, Telek, and Vassena
recently interpreted it as a layered mixed matrix and characterized the
nonzero pattern of its inverse in the independent-symbol model
\cite{kuhrs-telek-vassena-2026}.  Their Remark~3.6 also points out the precise
motivation here: classical mass action is not parameter-rich, so its
derivatives obey extra algebraic relations at a steady state.

We resolve a deliberately finite test of that distinction.  No source cited
above poses this six-complex census as an open problem, and we make no claim
about networks beyond its boundary.  From a certification perspective, the
method is akin to exact finite work in a different setting that combines
symmetry reduction, machine-readable certificates, and independent
reconstruction \cite{zeng-liu-ratnavelu-ong-2026}; here every decisive
predicate is likewise evaluated over the integers or rationals.

Index the ordered complex set by
\[
 c_0=(0,0),\quad c_1=(1,0),\quad c_2=(0,1),\quad
 c_3=(2,0),\quad c_4=(1,1),\quad c_5=(0,2).
\]
A reaction is a pair \((i,j)\), \(i\ne j\), representing
\(c_i\to c_j\), and reactions are ordered lexicographically.  A network
\(E\) is a three- or four-element subset of the 30 reactions.  The involution
\(\tau:X\leftrightarrow Y\) induces
\[
 (0,1,2,3,4,5)\longmapsto(0,2,1,5,4,3).
\]
We sort the reaction list and represent its \(\tau\)-orbit by the smaller of
\(E\) and \(\tau E\) in ordinary lexicographic order.  Canonical networks are
then ordered lexicographically, followed by their reactions and the species
order \(X<Y\).

For a sorted network with \(q\) reactions, let
\[
 (S^-)_{\bullet\rho}=c_{\mathrm{source}(\rho)},\qquad
 S_{\bullet\rho}=c_{\mathrm{target}(\rho)}-
 c_{\mathrm{source}(\rho)}.
\]
The generic reactivity matrix \(R_{\rm gen}\in\R^{q\times2}\) has an
independent symbol \(r_{\rho m}\) in position \((\rho,m)\) precisely when
species \(m\) occurs in the source of \(\rho\), and zero otherwise.  We retain
exactly the networks satisfying
\begin{equation}\label{eq:filters}
 \rank S=2,\qquad \ker S\cap\R_{>0}^{q}\ne\varnothing,
 \qquad \det(SR_{\rm gen})\not\equiv0.
\end{equation}
Choose a rational full-column-rank matrix
\(G\in\Q^{q\times(q-2)}\) with \(SG=0\), and set
\[
 A_{\rm gen}=[R_{\rm gen}\mid G].
\]
The first two rows of \(A_{\rm gen}^{-1}\), indexed by \(X,Y\), do not
depend on the chosen Gale basis as far as their zero pattern is concerned.

At a positive mass-action steady state, write \(v>0\) for the reaction-flux
vector and \(x=(x_X,x_Y)>0\) for the concentration vector.  Then
\begin{equation}\label{eq:ma}
 R_{\rm MA}(v,x)=\diag(v)(S^-)^T\diag(x)^{-1},
 \qquad Sv=0,
 \qquad A_{\rm MA}=[R_{\rm MA}\mid G].
\end{equation}
Conversely, every such \((v,x)\) is realized by positive rate constants:
for a reaction with source complex \(c_i\), take
\(k_\rho=v_\rho/x^{c_i}\).  Moreover,
\begin{equation}\label{eq:normalization}
 A_{\rm MA}(v,x)=A_{\rm MA}(v,(1,1))
 \diag(x_X^{-1},x_Y^{-1},I_{q-2}),
\end{equation}
so concentration normalization preserves zeros in the species rows of the
inverse.

\begin{theorem}\label{thm:main}
Among all networks retained by \eqref{eq:filters}, there is a pair for which
an entry of \(A_{\rm gen}^{-1}\) is nonzero but the corresponding entry of
\(A_{\rm MA}^{-1}\) vanishes identically on the positive steady-state
parameter variety.  In the ordering above, the lexicographically minimum
triple is
\[
 E_*=(0\to X,\ 0\to2X,\ X\to Y,\ X+Y\to0),
 \qquad \rho=0\to X,\qquad m=Y.
\]
The mass-action inverse is generically defined for this network; indeed, its
determinant is nonzero at every positive steady flux.
\end{theorem}

\section{The exact witness}

Label the displayed reactions of \(E_*\) by \(\rho_0,\ldots,\rho_3\).  Direct
construction gives
\begin{equation}\label{eq:matrices}
 S=\begin{pmatrix}1&2&-1&-1\\0&0&1&-1\end{pmatrix},\qquad
 S^-=\begin{pmatrix}0&0&1&1\\0&0&0&1\end{pmatrix},\qquad
 G=\begin{pmatrix}2&2\\-1&0\\0&1\\0&1\end{pmatrix}.
\end{equation}
The minor of columns 0 and 2 of \(S\) equals 1, while \(SG=0\).  At
\(z=(-1,2)^T\),
\[
 v=Gz=(2,1,2,2)^T>0,
\]
so the first two filters hold.

Let
\[
 a=r_{\rho_2X},\qquad b=r_{\rho_3X},\qquad c=r_{\rho_3Y}.
\]
These are exactly the nonzero independent symbols, and
\[
 R_{\rm gen}=\begin{pmatrix}0&0\\0&0\\a&0\\b&c\end{pmatrix},
 \qquad
 A_{\rm gen}=\begin{pmatrix}
 0&0&2&2\\
 0&0&-1&0\\
 a&0&0&1\\
 b&c&0&1
 \end{pmatrix}.
\]
Expansion yields
\begin{equation}\label{eq:genericdet}
 \det(SR_{\rm gen})=\det A_{\rm gen}=2ac.
\end{equation}
Thus \(E_*\) passes the last filter.  The signed adjugate cofactor for inverse
row \(Y\) and reaction column \(\rho_0\) is
\[
 C_{\rho_0,Y}=-a+b,
 \qquad
 (A_{\rm gen}^{-1})_{Y,\rho_0}=\frac{-a+b}{2ac},
\]
which is not the zero rational function in the independent symbols.

At \(x=(1,1)\), mass action specializes
\[
 a=v_2,\qquad b=v_3,\qquad c=v_3.
\]
The cofactor then satisfies the polynomial identity
\begin{equation}\label{eq:identity}
 C^{\rm MA}_{\rho_0,Y}=-v_2+v_3=-(Sv)_Y.
\end{equation}
It therefore vanishes on the entire linear steady-flux variety \(Sv=0\), not
merely at a sampled point.  On that variety, writing \(v=Gz\) gives
\begin{equation}\label{eq:madet}
 \det A_{\rm MA}(v,(1,1))=2z_1^2=2v_3^2.
\end{equation}
Every positive steady flux has \(v_3>0\), so the inverse is defined there.
Equation~\eqref{eq:normalization} extends both the zero and nondegeneracy
statements to every \(x>0\).  This proves the existential part of
Theorem~\ref{thm:main}.

For completeness, the only earlier pair within this same canonical network is
\((\rho_0,X)\).  Its generic cofactor is \(-c\), and under mass action it is
\(-v_3\), which never vanishes for a positive steady flux.  Thus it does not
precede the stated witness.

\section{Exact census and minimality}

There are
\[
 \binom{30}{3}+\binom{30}{4}=31{,}465
\]
labelled networks.  Exactly 133 are fixed by \(\tau\), so Burnside's lemma
gives \((31{,}465+133)/2=15{,}799\) canonical orbits.  A complete exact
enumeration sorted these representatives and applied the following tests.

\begin{lemma}\label{lem:tests}
For \(q\in\{3,4\}\), all retention, inverse-cofactor, and mass-action
vanishing predicates in the census reduce to finite rational arithmetic.
\end{lemma}

\begin{proof}
Rank and a Gale basis are obtained by rational row reduction.  When
\(q=3\), strict positivity of the one-dimensional kernel is decided by the
two orientations of a basis vector.  When \(q=4\), the inequalities
\(Gz>0\) define an open cone in \(\Q^2\).  If it is nonempty, it contains a
point with \(z_0\ne0\): this is immediate unless a chosen feasible point has
\(z_0=0\), in which case openness permits a small rational perturbation.
After scaling, it is therefore enough to intersect the inequalities with
\(z_0=1\) and \(z_0=-1\).  Each intersection is an exact open interval with
rational endpoints.

Writing \(S_{\bullet j}\) for a stoichiometric column, expansion gives
\[
 \det(SR_{\rm gen})=
 \sum_{j,k}\det(S_{\bullet j},S_{\bullet k})
 r_{jX}r_{kY},
\]
where only symbols allowed by the source complexes occur.  The symbols are
independent, so nonvanishing is decided by the integer coefficients.

After deletion of reaction row \(\rho\) and species column \(m\) from
\([R_{\rm gen}\mid G]\), only the other reactivity column remains.  Its
determinant is linear in distinct independent symbols, and its coefficients
are rational minors of \(G\).  Under mass action it becomes a linear form
\(\alpha^Tv\).  Restriction to \(Sv=0\), parameterized by \(v=Gz\), is the
linear form \(\alpha^TGz\); it vanishes identically exactly when all of those
rational coefficients are zero.  The retained positive cone is a nonempty
relatively open subset of \(\ker S\), so this is also exactly the criterion
for identical vanishing on the positive locus.  Finally, direct determinant expansion of
\([R_{\rm MA}(Gz,(1,1))\mid G]\) decides whether the denominator is the zero
polynomial.  These are all exact finite tests.
\end{proof}

The full census produced the following audit counts.
\begin{center}
\small
\begin{tabular}{@{}p{0.72\textwidth}r@{}}
\toprule
Predicate & Count\\
\midrule
Canonical orbits & 15,799\\
Rank-two orbits & 15,607\\
Rank two with a strictly positive kernel vector & 5,176\\
Retained orbits & 5,168\\
Retained orbits with identically zero mass-action determinant & 463\\
Generic-nonzero inverse pairs checked & 35,968\\
Mass-action-zero pairs with nonzero denominator & 1,432\\
\bottomrule
\end{tabular}
\end{center}
The first qualifying pair occurs at zero-based orbit index 303, namely the
304th representative, and is the pair in Theorem~\ref{thm:main}.  No
representative at indices 0 through 302 contains a qualifying pair.

The main implementation writes one machine-readable record for each of the
15,799 representatives, including the canonical reaction list, filter
decisions, rational Gale data, an exact positive flux when one exists, the
mass-action determinant polynomial, and every generic-nonzero cofactor normal
form.  A second implementation shares neither the census code nor its linear
algebra library.  Using only rational fractions, separate Gaussian
elimination, and direct determinant formulas, it reconstructs all canonical
representatives and reclassifies every representative through index 303.  It
again finds no earlier witness and verifies \eqref{eq:matrices}--\eqref{eq:madet}.
Together with Lemma~\ref{lem:tests}, this establishes the claimed
lexicographic minimality.

\section{Scope, reproducibility, and assistance}

The theorem is exhaustive only for the six ordered complexes above, three or
four distinct loop-free directed reactions, the species-swap equivalence, and
the three filters in \eqref{eq:filters}.  It neither classifies larger
networks nor asserts that every failure of parameter richness creates a zero.
Networks on which the specialized determinant vanishes identically are
recorded but are not counted as witnesses, because a generic mass-action
inverse is then unavailable.

The accompanying files include the full census program and JSON-lines
manifest, a summary record, and the separately written prefix verifier with
its output.  Running the two Python programs from the project root reproduces
the classification and the minimum.  The main census uses exact SymPy
rational arithmetic; the replay uses Python's standard-library
\texttt{fractions.Fraction}.  No floating-point comparison or random choice
enters a decisive predicate.

OpenAI Codex assisted with literature search, search-space formalization,
development and checking of the exact enumeration, and manuscript
preparation.  Language-model output is not mathematical evidence: the result
rests on the displayed identities and the released exact-arithmetic
certificates.  No proof assistant was used.

\bibliographystyle{plain}
\bibliography{references}

\end{document}
